\documentclass[15pt]{article}
\usepackage[top=1in, bottom=1in, left=1in, right=1in]{geometry}
\usepackage[utf8]{inputenc}
\usepackage[T1]{fontenc}
\usepackage{amsmath}
\usepackage{amssymb}
\usepackage{lmodern}
\usepackage{bm}
\usepackage{scrextend}
%\usepackage{showframe}
\usepackage{calc}
\usepackage{changepage}
\usepackage{mdframed}
\usepackage{dsfont}
\usepackage{enumitem}
\usepackage{pbox}
\usepackage{setspace}
\usepackage{stmaryrd}
\usepackage{setspace}
\usepackage[table]{xcolor}
\usepackage{multicol}
\usepackage{mathtools}
\usepackage{hanging}
\usepackage{longtable}
\usepackage{aligned-overset}
\usepackage{ntheorem}
\usepackage{mdframed}
\usepackage[theorems,breakable]{tcolorbox}%
\usepackage[backend=bibtex,style=alphabetic]{biblatex}
\addbibresource{lib.bib}

\allowdisplaybreaks
\hbadness=10000
\makeatletter
\renewcommand*\env@matrix[1][*\c@MaxMatrixCols c]{%
  \hskip -\arraycolsep
  \let\@ifnextchar\new@ifnextchar
  \array{#1}}

\newcommand{\mathleft}{\@fleqntrue\@mathmargin10pt}
\newcommand{\mathcenter}{\@fleqnfalse}
\makeatother

\newcommand*{\QED}{\hfill\ensuremath{\square}} 
\newcommand{\R}{\mathbb{R}}
\newcommand{\N}{\mathbb{N}}
\newcommand{\Z}{\mathbb{Z}}
\newcommand{\Q}{\mathbb{Q}}
\newcommand{\p}{\mathcal{P}}
\newcommand{\K}{\mathbb{K}}
\newcommand{\C}{\mathbb{C}}
\newcommand{\la}{\langle}
\newcommand{\ra}{\rangle}
\newcommand{\E}{\mathbb{E}}
\newcommand{\1}{\mathds{1}}


\newcommand{\bol}[1]{\mathbf{#1}}


\newlength{\hangwidth}
\newcommand{\newhang}[1]{\settowidth{\hangwidth}{#1}\par\hangpara{\hangwidth}{1}#1}

\newcommand{\norm}[1]{\lVert #1 \rVert}

\def\Abstand{1cm}

\newtcbtheorem[auto counter]{lemma}{Lemma}{%
  theorem name and number,%
        colframe=gray,%
        fonttitle=\bfseries,
        boxsep=5pt,
        left=5pt, right=5pt, top=5pt, bottom=5pt
    }{lem}

\newtcbtheorem[auto counter]{theorem}{Theorem}{
  theorem name and number,%
        colframe=gray,%
        fonttitle=\bfseries,
        boxsep=5pt,
        left=5pt, right=5pt, top=5pt, bottom=5pt
    }{thm}

\newtcbtheorem[auto counter]{definition}{Definition}{
  theorem name and number,%
      colframe=gray,%
      fonttitle=\bfseries,
      boxsep=5pt,
      left=5pt, right=5pt, top=5pt, bottom=5pt
  }{def}

\title{Notes on Distributional Reinforcement Learning}
\author{ }
\date{}


\begin{document}


\maketitle

\newlength\breite
\setlength\breite{\linewidth-4pt}
\setlength\fboxsep{0pt}
\setlength\fboxrule{0.25pt}
\setlength{\abovedisplayskip}{3mm} %Abstand OBEN zw. Text und Formel
\setlength{\belowdisplayskip}{3mm} %Abstand OBEN zw. Text und Formel
\setlength\itemsep{0pt}
\setlength\parindent{0pt}


\mathleft

\section*{Preliminaries}

These notes follow the modern treatment of distributional reinforcement learning
in \cite{bellemare2023distributional}.

\section{Linear Function Approximation}

Let $(S_t,A_t,R_{t+1},S_{t+1})_{t\geq 0}$ be the trajectory of a Markov decision
process under a fixed policy $\pi$. For a discount factor $\gamma \in [0,1)$, the
return is
\[
  G_t = \sum_{k=0}^{\infty} \gamma^k R_{t+k+1}.
\]
The classical value function is the conditional mean of this random return:
\[
  v^\pi(s) = \E_\pi[G_t \mid S_t=s].
\]
Equivalently, it is the fixed point of the Bellman operator
\[
  (T^\pi v)(s)
  =
  \E_\pi\left[R_{t+1}+\gamma v(S_{t+1}) \mid S_t=s\right].
\]

In small finite state spaces one may store one value per state. In large or
continuous state spaces this is impossible, so we approximate the value function
from a feature map
\[
  \phi:\mathcal{S}\to \R^d,
  \qquad
  \phi(s) =
  \begin{pmatrix}
    \phi_1(s) \\ \vdots \\ \phi_d(s)
  \end{pmatrix}.
\]
A linear value approximation has the form
\[
  v_w(s) = \phi(s)^\top w,
  \qquad w\in\R^d.
\]
The parameter vector $w$ is shared across states. Thus learning at one state can
generalize to other states with similar features.

The temporal-difference error for a transition $(S_t,R_{t+1},S_{t+1})$ is
\[
  \delta_t
  =
  R_{t+1} + \gamma \phi(S_{t+1})^\top w_t
  - \phi(S_t)^\top w_t.
\]
The semi-gradient TD(0) update is
\[
  w_{t+1}
  =
  w_t + \alpha_t \delta_t \phi(S_t),
\]
where $\alpha_t>0$ is a step size. This update moves $v_{w_t}(S_t)$ toward the
one-step bootstrap target
\[
  R_{t+1} + \gamma v_{w_t}(S_{t+1}).
\]

The important point is that linear TD does not learn the full table
$s\mapsto v(s)$. It learns a point in the linear function class
\[
  \mathcal{V}_\phi
  =
  \{s\mapsto \phi(s)^\top w : w\in\R^d\}.
\]
When the true value function is not in this class, the algorithm seeks a
projected Bellman fixed point: roughly, the Bellman update is followed by a
projection back onto the functions representable by the features. The quality of
the approximation therefore depends strongly on the representation $\phi$.

\section{Distributional Value Approximation}

Classical reinforcement learning compresses the return distribution to its mean.
Distributional reinforcement learning keeps the random return itself. For a fixed
policy $\pi$, define
\[
  Z^\pi(s) \stackrel{D}{=} G_t \mid S_t=s,
\]
where $\stackrel{D}{=}$ denotes equality in distribution. Its expectation
recovers the classical value:
\[
  v^\pi(s) = \E[Z^\pi(s)].
\]
The distributional Bellman operator is
\[
  (T^\pi Z)(s)
  \stackrel{D}{=}
  R_{t+1} + \gamma Z(S_{t+1}),
  \qquad S_t=s.
\]
Thus the Bellman target is no longer a scalar. It is the distribution of a random
reward plus a discounted next-state return distribution.

To approximate distributions, one common choice is to represent each return
distribution by its quantile function. For a real-valued random variable $Z$ with
cumulative distribution function $F_Z$, its quantile function is
\[
  F_Z^{-1}(\tau)
  =
  \inf\{z\in\R : F_Z(z)\geq \tau\},
  \qquad \tau\in(0,1).
\]
Choose quantile levels
\[
  0 < \tau_1 < \cdots < \tau_N < 1.
\]
A quantile approximation represents the return distribution at state $s$ by
\[
  Z_\theta(s)
  =
  \frac{1}{N}\sum_{i=1}^N \delta_{\theta_i(s)},
\]
where $\delta_x$ is a Dirac mass at $x$, and $\theta_i(s)$ is intended to
approximate $F^{-1}_{Z^\pi(s)}(\tau_i)$.

With linear function approximation, each quantile location can be parameterized
linearly:
\[
  \theta_i(s) = \phi(s)^\top w_i,
  \qquad w_i\in\R^d.
\]
Equivalently, the parameters are
\[
  W = (w_1,\ldots,w_N)\in\R^{d\times N}.
\]
This is the distributional analogue of linear value approximation: instead of
one linear prediction for the mean return, we use several linear predictions for
different parts of the return distribution.

\section{Quantile TD}

Quantile TD is the temporal-difference idea applied to quantile
representations. Given a transition
\[
  (S_t,R_{t+1},S_{t+1}),
\]
the one-step target samples are
\[
  y_j
  =
  R_{t+1} + \gamma \theta_j(S_{t+1}),
  \qquad j=1,\ldots,N.
\]
For the current state, the predictions are
\[
  x_i = \theta_i(S_t),
  \qquad i=1,\ldots,N.
\]
The quantile regression loss for quantile level $\tau$ is
\[
  \rho_\tau(u)
  =
  u\bigl(\tau-\1_{\{u<0\}}\bigr).
\]
For one transition, the distributional TD loss is commonly written as
\[
  L(W)
  =
  \frac{1}{N}
  \sum_{i=1}^N
  \sum_{j=1}^N
  \rho_{\tau_i}(y_j - x_i).
\]
The residual $y_j-x_i$ compares a target quantile sample from the next-state
distribution with a predicted quantile at the current state. Minimizing this
loss moves the predicted quantiles of $Z_\theta(S_t)$ toward the Bellman-updated
distribution
\[
  R_{t+1}+\gamma Z_\theta(S_{t+1}).
\]

For linear quantiles, a subgradient update has the form
\[
  w_{i,t+1}
  =
  w_{i,t}
  +
  \alpha_t
  \left(
    \frac{1}{N}\sum_{j=1}^N
    \bigl(\tau_i-\1_{\{y_j < x_i\}}\bigr)
  \right)
  \phi(S_t),
  \qquad i=1,\ldots,N.
\]
This mirrors ordinary TD:
\[
  \text{scalar TD error}
  \quad \leadsto \quad
  \text{quantile regression residuals}.
\]
The scalar target $R+\gamma v(S^{\prime})$ is replaced by a target distribution
$R+\gamma Z(S^{\prime})$, and the mean-squared TD objective is replaced by a quantile
regression objective.

In practice, the non-smooth loss $\rho_\tau$ is often replaced by a quantile
Huber loss, which behaves quadratically near zero and linearly for large
residuals. This improves numerical stability while preserving the asymmetric
penalty that makes quantile regression estimate quantiles rather than means.

\section{Gateway to Deep Representations}

The linear setting separates two ideas:
\[
  \text{features } \phi(s)
  \qquad\text{and}\qquad
  \text{linear prediction heads } w_i.
\]
Deep distributional reinforcement learning keeps this structure but learns the
features. A neural network first maps the state to a learned representation
\[
  h_\eta(s)\in\R^d,
\]
and then predicts either atom probabilities, quantile locations, or a sampled
quantile function from that representation. In the quantile case one may write
\[
  \theta_i(s) = f_i(h_\eta(s)).
\]

This is the conceptual bridge to deep representations: the question is no longer
only how to update a fixed set of quantiles, but also what representation makes
return distributions easy to predict. The linear model asks, ``Which return
distributions are expressible using these features?'' The deep model asks,
``Can the agent learn features for which distributional prediction becomes
simple?''

\printbibliography

\end{document}
